\answer{ex:14:1}{(a)~$0.17$~s と $1.5$~s。(b)~$11$~cm 未満の距離から。}
\answer{ex:14:2}{$g\le\beta w_{q\mu}/w_{z\mu}=5$ である限り、痛みのない触覚はどんな $\mu$ でも通らない。触覚 $\mu=1$ は $g>6$ で通る。強い感作はふつうの接触を痛みに変える。}
\answer{ex:14:3}{(a)~$g^*=(1-k_sC)/(1-k_sA)$、$k_sA<1$ で安定。(b)~触覚なしでは $\nu\ge2$、触覚ありではすでに $\nu\ge1.7$ で。}
\answer{ex:14:4}{$V(t)=-Pe^{-(T-t)/\tau_\gamma}$。(a)~$-Pe^{-T/\tau_\gamma}\approx-0.37P$。(b)~$+Pe^{-(T-t_e)/\tau_\gamma}\approx+0.61P$。これは $t_e$ とともに $+P$ まで増え、痛みを避けることを教える。}
\answer{ex:14:5}{$k_s=4$ ではラッチしない（第2の安定状態は $k_s\ge4.58$ で現れる）。$k_s=4.6$, $5$, $6$, $8$ で $T_{\min}\approx4.35$, $1.40$, $0.65$, $0.32\,\tau_s$。}
\answer{ex:14:6}{$w_{q\nu}=4/9\approx0.444$, $q_0=2/9\approx0.222$, $w_{q\mu}=49/45\approx1.089$。痛みのない触覚は $g=49/9\approx5.4$ まで安全。}
\answer{ex:14:7}{(a)~しきい値は $\nu_c\ge q_0/w_{q\nu}$ のとき $\nu_c(g)=\theta/g$、そうでなければ $(\theta+\beta q_0)/(g+\beta w_{q\nu})$。$g^*\approx2.45$, $1.0$, $0.25$。(b)~未解決の問題。}
